
    \Quote{Trabajar en calidad significa diseñar, producir y servir un producto o servicio que sea útil, lo más económico posible y siempre satisfactorio para el usuario}{Kaoru Ishikawa}

    \resumebox{definicion}{Calidad}{Conjunto de valores y hábitos que, complementarios con el uso de prácticas y herramientas de calidad, permite el correcto cumplimiento de metas y objetivos}

\section{Costos de la calidad}

Hay varios tipos de costos relacionados a la calidad. Principalmente se pueden dividir en los costos por tener calidad y los costos por NO tener calidad.

En el caso de los primeros podemos encontrar costos por fallas tanto internas como externas. En el caso de las primeras podemos encontrar el reproceso de productos, el desperdicio de estos o los tiempos ociosos. En el caso de las segundas encontramos el costo de las devoluciones y reclamaciones de productos, el pago de garantías, así como las rebajas.

Por otro lado, los gastos por tener calidad están principalmente en la prevención y la evaluación. La prevención incluye cosas como la planeación de estrategias de caldidad, la capacitación del personal, la revisión de los productos y las auditorías (obtención y análisis de datos) La evaluación consisite en la verificación de los proveedores e inspecciones locales.